\subsection{Retrospective adaptation exposes a target--source trade-off}
\label{sec:adaptation-followup}

We next studied adaptation on a separate Kim library using fixed intended-allele
and experimental-context groups. This is a retrospective development experiment,
not an independently collected confirmation study. Its outer validation contains
5{,}561 groups and 23{,}120 distinct candidates, including all-zero groups; it is
distinct from both the original benchmark and the previously inspected adaptation
test. Checkpoint selection uses validation groups \emph{inside} the adaptation
label budget, while the exposed outer validation informs research-level decisions.
For the initial acquisition, nominal budgets of 200 and 1{,}000 groups contain
202 and 1{,}001 groups,
respectively: 168/34 and 839/162 training/inner-validation groups, requiring 788
and 3{,}848 distinct labelled candidates. These are per-fit adaptation costs, not
the total label cost of developing the method.

This low-label question is distinct from the earlier full-budget adaptation
result. With 19{,}357 labelled training groups and 5{,}561 labelled validation
groups, ordinary fine-tuning achieved mean top-choice efficiency 0.04178 across
three optimizer seeds versus released OptiPrime's 0.04020 on the separate
5{,}557-group adaptation test; each seed's paired interval was positive. That
retrospective, target-supervised improvement remains valid. It does not establish
low-label superiority, matched-supervision architectural superiority, or
independent cross-study transfer, and its test efficiencies must not be compared
directly with the different v3 validation population below.

The v3 screen comprised 34 fits, all with 100 target updates. Four arms were
replicated across three optimizer seeds at both budgets, using the same acquisition
subset. The source-anchored selector adds a zero-initialized residual readout to a
fixed positive affine transformation of our frozen source score. Its backbone is
frozen; the same final score selects candidates in both domains. No OptiPrime
predictions, features or parameters are used during adaptation. Source replay and
retention validation exclude components connected to source fold 0 or target
alleles/full spacers; retention-validation examples were seen in pretraining.

At 1{,}000 nominal groups, mean achieved top-choice efficiency over optimizer seeds
is 0.039754 for ordinary ordinal fine-tuning (A), 0.041465 for a shared
Huber-plus-ranking score (D), 0.041199 for the anchored pairwise selector (E), and
0.041254 for E with soft pairwise source distillation (F). Released OptiPrime reaches
0.041409 on these identical candidates. E improves on A by 0.001445
(95\% paired component-bootstrap interval [0.000975, 0.001950]), but its difference
from OptiPrime is $-0.000210$ [${-}0.000759$, 0.000327]. A nonsignificant difference
does not establish equivalence. The released comparator receives no target
adaptation labels, so this comparison is not matched for target supervision.

Source retention is evaluated on the actually deployed selector using 1{,}463
informative source-audit groups. E reaches 0.132948; the original checkpoint
reaches 0.133635. The change of $-0.000687$ has interval
[${-}0.002613$, 0.001337] and does not establish noninferiority at the working
$-0.001$ margin. F loses a further 0.003896 relative to E
([${-}0.006337$, ${-}0.001795$]) despite reducing source-replay pairwise KL in a
first-seed diagnostic. D has greater source loss, reaching 0.112274. These
findings motivate directly testing preservation of useful choices, rather than
treating agreement in soft pairwise probabilities as evidence of retained utility.
They do not identify a unique mechanism: source-domain mixture, loss scale and
checkpoint selection remain alternatives. Indeed, E trails A on the separate
source-retention validation mixture, despite leading it on the source audit.

The replicated values above are averages of achieved outcomes, not a score
ensemble. A separate equal-weight E score ensemble reaches target efficiency
0.041312, also without a reliable OptiPrime advantage. All intervals in this
development analysis are exploratory and unadjusted. The archived v3 protocol,
run-level results and subsequent optimization study are versioned separately
from the original benchmark artifacts.

\subsubsection{Stronger optimization controls narrow the architectural claim}

The subsequent v3.1 screen crossed five methods with two learning-rate
multipliers and two horizons (100 or 500 target updates), yielding 20 fits on
the first 1{,}001-group acquisition and one optimizer seed. Each trajectory
offered the original deployment plus ten evenly spaced checkpoint choices;
hyperparameters were selected by inner target utility, not outer validation.
Ordinary fine-tuning A then reached target efficiency 0.041101, close to E's
0.041304. Their difference was 0.000203
([${-}0.000295$, 0.000736]), so the target advantage observed under v3's fixed
optimization budget was not established after this search. E retained more
source-audit utility than A (difference 0.004746 [0.002327, 0.007184]), but again
trailed it on the different source-validation mixture. The fresh frozen head
had the highest conventional inner utility, yet its selected outer efficiency
was only 0.038883. These results expose model-selection variability rather than
supporting promotion of a post-hoc outer-validation winner.

We next held E's target recipe fixed and tested twelve combinations of source
objective (soft pairwise KL, teacher top-choice margin, or label-aware expected
regret), study sampling (mixture or balanced), and penalty weight (0.1 or 1).
The teacher was our own frozen source checkpoint. A constrained checkpoint
policy required source-validation point loss no greater than 0.001 in each of
two studies, then maximized inner target utility, with original deployment as
fallback. Balanced teacher-margin preservation at weight 1 was the only recipe
to select a nonzero feasible checkpoint in this screen. This point-estimate
filter is not a statistical noninferiority guarantee, and equal scalar weights
did not match gradient pressure across objectives. The resulting recipe was
frozen before acquisition replication; it is not an OptiPrime hybrid.

\subsubsection{Acquisition replication reveals a preservation trade-off}

The locked replication crossed two total label budgets, three acquisition
subsets and two optimizer seeds for A, the inner-selected fresh frozen control,
E and balanced margin preservation. Four exact screen configurations were
reused, requiring 44 new fits and bringing v3.1 to 76 fits in total. The
acquisitions contain 200--202 or 1{,}000--1{,}002 groups, including inner
validation. Their union comprises 2{,}853 groups and 10{,}507 labelled
candidates; these development costs exclude the additional 23{,}120 outer
validation candidates and all source/pretraining labels.

At the 1{,}000-group budget, target-only selection gives mean target efficiency
0.041233 for tuned A, 0.041047 for E and 0.041126 for margin preservation
(Table~\ref{tab:v31-replication}). E's difference from A is now
$-0.000186$ [${-}0.000506$, 0.000157]. Its source advantage also fails to
persist: source efficiency averages 0.128462 for E and 0.128679 for A, versus
initialization at 0.133635. Across acquisition subsets, E's source means are
0.133997, 0.117906 and 0.133483. Thus a frozen backbone and an anchored score do
not alone guarantee retention of useful source choices.

Margin preservation improves source efficiency relative to E by 0.004822
[0.003636, 0.006004] under target-only selection, with target difference
0.000079 [${-}0.000102$, 0.000266]. Its own source change relative to
initialization is $-0.000351$ [${-}0.001626$, 0.001106], so this policy does not
clear the working retention margin. Its target difference from released
OptiPrime is $-0.000284$ [${-}0.000809$, 0.000237], which does not establish
superiority or equivalence. Pooled target Spearman remains higher for A and
margin preservation (0.732025 and 0.712110) than OptiPrime (0.695996), illustrating
again that better pooled prediction is not sufficient evidence of better @1.

With source-constrained selection, margin preservation adapts in five of six
1{,}000-group runs, whereas constrained A, fresh frozen and E always fall back.
The margin procedure reaches target/source efficiencies 0.040652/0.133858.
Its averaged source-audit change is 0.000223 [${-}0.000821$, 0.001351], whose
lower bound exceeds the working $-0.001$ margin on this exposed audit.
However, target utility falls below tuned A by 0.000580
([${-}0.000883$, ${-}0.000264$] for margin minus A), and below OptiPrime by
0.000757 ([${-}0.001324$, ${-}0.000219$]). At budget 200, the source interval
does not clear the retention margin. This is a measurable trade-off, not a
joint target-superiority and source-retention result.

All replication intervals condition on the six fitted runs and resample the
same evaluation components after averaging their outcomes; they do not capture
all possible acquisitions or constitute independent-study confirmation. The
stronger RNA/DNA-normalized 19-base-spacer overlap audit identifies one additional
target component (one group, eight candidates); excluding that entire component
does not change the qualitative comparisons. No matching source row enters
informative replay. Independent public-data confirmation remains gated by
linked outcome availability and input compatibility, so larger-budget and
architecture branches were deferred, not counted as failed experiments.

\begin{table}[htbp]
\centering
\small
\caption{Locked-recipe acquisition replication on the exposed external-library validation population. Each entry averages achieved outcomes over three acquisition subsets and two optimizer seeds, not averaged scores. All efficiencies are fractions. Target-only (T) and per-study source-constrained (C) checkpoint policies are distinct procedures. Fallback counts indicate deployment of the unchanged starting model.}
\label{tab:v31-replication}
\begin{tabular}{rlcrrrr}
\toprule
Budget & Method & Policy & Target @1 & $\Delta$ OptiPrime & Source @1 & Fallbacks \\
\midrule
200 & Ordinal A & T & 0.040158 & $-0.001251$ & 0.127608 & 0/6 \\
200 & Fresh frozen & T & 0.039208 & $-0.002201$ & 0.129882 & 4/6 \\
200 & Anchored E & T & 0.039381 & $-0.002028$ & 0.127029 & 1/6 \\
200 & Margin & T & 0.039669 & $-0.001740$ & 0.132864 & 2/6 \\
200 & Margin & C & 0.039853 & $-0.001556$ & 0.133149 & 2/6 \\
\midrule
1,000 & Ordinal A & T & 0.041233 & $-0.000177$ & 0.128679 & 0/6 \\
1,000 & Fresh frozen & T & 0.039800 & $-0.001609$ & 0.117924 & 0/6 \\
1,000 & Anchored E & T & 0.041047 & $-0.000363$ & 0.128462 & 0/6 \\
1,000 & Margin & T & 0.041126 & $-0.000284$ & 0.133284 & 0/6 \\
1,000 & Margin & C & 0.040652 & $-0.000757$ & 0.133858 & 1/6 \\
\bottomrule
\end{tabular}
\par\smallskip
\begin{minipage}{0.98\textwidth}\footnotesize
Released OptiPrime target @1 is 0.041409; starting-model target/source @1 are 0.039627/0.133635.
Target evaluation retains all 5,561 groups; source audit uses 1,463 informative groups. All checkpoint policies, run-level outcomes, conditional component intervals and acquisition-specific results are reported in the v3.1 replication ledger. Released OptiPrime receives no target adaptation labels.
\end{minipage}
\end{table}

